\section{Provenance: Is the Payload Hostile?}
\label{sec:malware}

\subsection{The hypothesis on arrival}

The files came from a third party's Windows machine. The stated concern was a
virus, and the specific theory was that the third party had downloaded from an
impersonation site---a lookalike domain serving a reskinned payload under the
name of a known repacking group. That theory has a specific and testable
consequence: an impersonated repack would not match the genuine distributor's
published checksums, and would likely contain persistence mechanisms absent from
the genuine article.

A secondary theory, more sophisticated and more alarming, was raised early: that
the repacking process might involve \emph{re-implementing} a commercial
anti-tamper system, and that such a re-implementation would be an ideal vehicle
for injecting code at runtime. This theory is the reason the anti-tamper question
(Section~\ref{sec:denuvo}) had to be answered even though it initially looked
like a pure compatibility concern.

\subsection{Method}

Analysis was static and offline. The full triage comprised six independent
lines of evidence.

\subsubsection{Checksum provenance}
All six containers were verified against the distributor's published MD5
manifest. All six matched. This does not establish that the distributor is
benign---it establishes that the files are the ones the distributor published,
which directly falsifies the impersonation-site theory.

\subsubsection{Installer script recovery}
The installer is an Inno Setup~\cite{innosetup} package, whose compiled script is
not visible to \verb|strings| and is not the first compressed object in the file.
The recovery procedure is reusable and is documented here because the obvious
approach fails:

\begin{quote}
The \verb|zlb\x1a| magic near the start of the overlay marks the \emph{file
data}, not the script. Locate the SetupLdr offset table by searching for the
12-byte identifier \verb|rDlPtS|, then read nine little-endian 32-bit words:
version, total size, EXE offset, uncompressed EXE size, EXE CRC,
\textbf{offset0} (script header), offset1 (file data), and a table CRC. At
offset0, skip the 64-byte \verb|Inno Setup Setup Data (x.y.z) (u)| banner, then
read a 4-byte CRC, a 4-byte stored size, and a 1-byte compression flag, followed
by chunks of \verb|[crc32(4) + up to 4096 bytes]| until the stored size is
satisfied. The concatenated payload is raw LZMA1 whose first five bytes are the
properties block.
\end{quote}

The decompressed script is the authoritative statement of everything the
installer will do to the system, and reading it converts the persistence question
from a dynamic-analysis problem into a text-search problem.

\subsubsection{Persistence audit}
The recovered script contains exactly one registry write:
\verb|AppCompatFlags\Layers| set to \verb|RUNASADMIN|. There are no \verb|Run|
keys, no \verb|Winlogon| modifications, no service installations, and no
scheduled tasks.

An important negative: string searches for \verb|Startup| and \verb|RegSvr32| do
hit, and in an automated report those hits would be flagged. They are matches
against Inno Setup's own built-in English message table---the strings the
installer would display if it were performing such an action---not evidence that
it performs one. \textbf{A string table is not a call site.} This distinction
accounted for the majority of the initial false positives.

\subsubsection{Network endpoint enumeration}
Every URL in the compiled script resolves to one of three categories: Microsoft
redistributables (Visual C++ runtimes, .NET), NVIDIA PhysX, or the distributor's
own domain. The Inno Download Plugin is present and wired up---so the installer
genuinely can fetch---but the only things it is instructed to fetch are those
redistributables.

\subsubsection{The alarming-but-benign component}
Two files, \verb|host.cmd| and \verb|hosts.exe|, are the sort of thing that
correctly triggers alarm: a program that edits the system \verb|hosts| file.
Reading them settles it. They point roughly sixty lookalike domains at a single
address and then \emph{remove} any existing entry that blocks the distributor's
real domain. This is an anti-impersonation measure authored by the distributor.

The finding is doubly useful. It is benign, and its presence is positive evidence
that the repack is the genuine article---which is exactly the question
Section~\ref{sec:malware} opened with. \textbf{The most alarming artefact in the
package turned out to be the strongest disproof of the theory that motivated the
investigation.}

\subsubsection{Differential comparison against a known-clean package}
Twenty-eight leftover temporary directories were found in the Wine prefix, all
dated to a single day, indicating the installer had been run and had failed
roughly nine times. Every executable and library in those directories was
compared byte-for-byte against the payload of an independently verified clean
repack from the same distributor. All matched, including the archive
decompressor, the download plugin, and the hosts utility.

Four files differed, and each difference is explicable: three are per-repack
configuration (one of which contains the author's own installation path, so it
\emph{must} differ), and the fourth is a newer build of a compression tool with
added codec support.

\subsection{An anomaly correctly excluded}

An earlier pass had flagged a \verb|Winlogon| shell value set to an empty
string---a genuine and serious-looking anomaly. Timestamp comparison excluded it:
the modification predates the arrival of these files by twelve days. It is
unrelated.

We record this because the temptation, when investigating a suspected compromise,
is to attribute every anomaly to the incident under investigation. A system of
any age contains unrelated anomalies. \textbf{Establishing that an artefact is
older than the incident is a complete defence, and it is usually one timestamp
away.}

\subsection{Verdict}

\begin{table}[htbp]
\caption{Malware triage: six independent lines, all negative}
\label{tab:malware}
\centering
\footnotesize
\begin{tabularx}{\columnwidth}{@{}L{2.4cm}X@{}}
\toprule
\textbf{Line of evidence} & \textbf{Result} \\
\midrule
Checksum provenance & All six containers match the distributor's published
                      manifest; impersonation theory falsified \\
Installer script & Recovered in full; one registry write, and it is a
                   compatibility flag \\
Persistence & None. No Run keys, services, scheduled tasks, or Winlogon
              modification \\
Network & All endpoints are Microsoft, NVIDIA, or the distributor \\
\verb|hosts| tooling & Anti-impersonation, authored by the distributor;
                        evidence \emph{for} authenticity \\
Byte comparison & Payload identical to an independently verified clean package,
                  modulo per-repack configuration \\
Runtime footprint & Touches only its own directory, system libraries, and its own
                    application-data folder \\
\bottomrule
\end{tabularx}
\end{table}

\textbf{The files are clean.} No malware, no persistence, no unexpected network
behaviour. The investigation is closed and should not be reopened without new
evidence.

The remaining theory---that the repack re-implements an anti-tamper system, and
that the re-implementation is the injection vector---is not addressed by any of
the above. It is addressed, and destroyed, in the next section, by a single
observation that also happens to answer the compatibility question.
